\documentclass[11pt,a4paper]{article}
\usepackage[utf8]{utf8}
\usepackage[margin=0.9in]{geometry}
\usepackage{amsmath,amsfonts,amssymb}
\usepackage{graphicx}
\usepackage{hyperref}
\usepackage{listings}
\usepackage{booktabs}
\usepackage{xcolor}
\usepackage{cite}
\usepackage{setspace}

\setstretch{1.15}

\hypersetup{
    colorlinks=true,
    linkcolor=blue,
    citecolor=blue,
    urlcolor=blue
}

\lstset{
    basicstyle=\ttfamily\small,
    breaklines=true,
    frame=single,
    keywordstyle=\color{blue},
    commentstyle=\color{green!60!black},
    stringstyle=\color{orange}
}

\title{\Large\bfseries blackbox-sec: Parallel Derivative-Free Optimization for Automated Red-Teaming and Adversarial Robustness Testing of Black-Box LLMs}

\author{
    \textbf{Kamran Saberifard}\\[4pt]
    Department of Computer Engineering (Cybersecurity), MehrAstan University, Guilan, Iran\\
    Aryorithm Group — Global AI Infrastructure \& Security Research\\
    \href{mailto:kamisaberi@gmail.com}{kamisaberi@gmail.com} \quad|\quad ORCID: \href{https://orcid.org/0009-0002-7822-6168}{0009-0002-7822-6168}
}

\date{\today}

\begin{document}

\maketitle

\begin{abstract}
\noindent Commercial Large Language Models (LLMs) hosted behind REST APIs (e.g., OpenAI GPT-4o, Anthropic Claude 3.5, Google Gemini 1.5) do not expose internal neural network weights or log-probability gradients ($\nabla_x L$). Consequently, white-box gradient-based adversarial attack methods (e.g., Greedy Coordinate Gradient - GCG) are completely inapplicable for security red-teaming on proprietary cloud endpoints. Security auditing of commercial LLMs requires automated, query-efficient derivative-free optimization frameworks capable of navigating discrete, high-dimensional prompt spaces. In this paper, we introduce \textbf{blackbox-sec}, an open-source, high-performance C++/CUDA framework for parallel black-box adversarial red-teaming. \texttt{blackbox-sec} combines high-dimensional Trust Region Bayesian Optimization (TuRBO), Covariance Matrix Adaptation Evolution Strategy (CMA-ES), and Bayesian Optimization and Hyperband (BOHB) early stopping to systematically explore non-differentiable prompt perturbation landscapes. The engine maps continuous latent vectors $x \in [-1, 1]^d$ to discrete token mutations (homoglyph insertions, zero-width space additions, and role-play framing) and offloads candidate response evaluation to custom CUDA kernels on GPU hardware. Empirical benchmarks demonstrate that \texttt{blackbox-sec} achieves a 94.8\% jailbreak discovery success rate on commercial LLMs using **86.7\% fewer API queries** than standard random search techniques, completing multi-trust-region campaigns with sub-millisecond C++ search worker overhead.
\end{abstract}

\vspace{6pt}
\noindent\textbf{Keywords:} Adversarial Machine Learning, Red-Teaming Automation, Derivative-Free Optimization, TuRBO, BOHB, CMA-ES, CUDA Acceleration, LLM Robustness, Black-Box Security.

\hr

\section{Introduction}
As Large Language Models (LLMs) are integrated into mission-critical corporate infrastructure, evaluating their robustness against adversarial manipulation has become a paramount security priority. Red-teaming refers to the practice of systematically probing AI models to identify safety boundary failures, jailbreak vulnerabilities, and unexpected behavior.

In white-box settings (where model weights are accessible), security researchers utilize gradient-based optimization algorithms like Greedy Coordinate Gradient (GCG)~\cite{gcg_attack} to compute exact token loss gradients and construct adversarial suffixes. However, enterprise applications overwhelmingly consume proprietary black-box LLMs hosted behind commercial REST APIs. These endpoints withhold model weights, internal activations, and token-level log-probabilities.

\begin{figure}[h!]
\centering
\fbox{\parbox{0.88\linewidth}{
\small
\textbf{[ Base Adversarial Prompt / Security Query ]}\\
$\Downarrow$ \textit{1. TuRBO Local Trust Region Sampling (turbo\_optimizer.cpp)}\\
\textbf{[ Continuous Latent Vector $x \in [-1, 1]^d$ ]}\\
$\Downarrow$ \textit{2. Derivative-Free Prompt Mutator (prompt\_mutator.cpp)}\\
\textbf{[ Perturbed Candidate Prompts (Homoglyphs + Framing) ]}\\
$\Downarrow$ \textit{3. Parallel API Queries to Target Black-Box LLM}\\
\textbf{[ Candidate Responses Evaluated on GPU CUDA Cores ]}\\
$\Downarrow$ \textit{4. Objective Fitness $F(x)$ Feedback $\rightarrow$ TuRBO Contraction/Expansion}\\
\textbf{\textcolor{red}{[ DISCOVERED JAILBREAK PAYLOAD & EXECUTIVE REPORT ]}}
}}
\caption{\texttt{blackbox-sec} Parallel Derivative-Free Red-Teaming Pipeline.}
\label{fig:blackbox_pipeline}
\end{figure}

\subsection{The Black-Box Red-Teaming Bottleneck}
Existing black-box red-teaming approaches suffer from two primary limitations:
\begin{enumerate}
    \item \textbf{Curse of High Dimensionality:} Simple random search or genetic algorithms fail to scale as prompt lengths and token vocabulary sizes grow. The search space expands exponentially, causing optimization stagnation.
    \item \textbf{Prohibitive Query Costs:} Querying commercial LLM APIs thousands of times per red-teaming campaign incurs significant financial costs and triggers API rate limits.
\end{enumerate}

To overcome these challenges, we present \textbf{blackbox-sec}, a C++20 and CUDA framework that unifies high-dimensional **Trust Region Bayesian Optimization (TuRBO)** with **BOHB successive halving** to automate black-box red-teaming with minimal query overhead.

\section{Mathematical Formulation \& Optimization Engine}

We formulate black-box adversarial red-teaming as a non-convex, non-differentiable continuous-to-discrete optimization problem:
\begin{equation}
\max_{x \in [-1, 1]^d} F\Big( \text{Mutate}(x_{\text{base}}, x) \Big) \quad \text{s.t.} \quad D(x, x_{\text{orig}}) \le \epsilon
\end{equation}
where $F(x) \in [0.0, 1.0]$ is an objective fitness function measuring safety policy violation severity, $\text{Mutate}(x_{\text{base}}, x)$ maps continuous point $x$ to discrete prompt perturbations, and $D(x, x_{\text{orig}})$ enforces semantic similarity constraints.

\subsection{1. High-Dimensional TuRBO Search}
Standard Gaussian Process (GP) surrogates fail in high dimensions due to $O(N^3)$ matrix inversion complexity and over-exploration of unpromising regions. \texttt{blackbox-sec} implements **TuRBO-m** (Trust Region Bayesian Optimization with $m$ local trust regions)~\cite{turbo_paper}.

Each trust region $TR_k$ maintains a local hyper-cube centered at its best observed point $x_k^{(b)}$ with dynamic side-length $L_k$. Candidate points are sampled strictly inside $TR_k$:
\begin{equation}
TR_k = \left[ x_k^{(b)} - \frac{L_k}{2}, x_k^{(b)} + \frac{L_k}{2} \right] \cap [-1, 1]^d
\end{equation}

The trust region length $L_k$ expands or contracts based on consecutive success/failure counters:
\begin{equation}
L_k^{(t+1)} = 
\begin{cases} 
\min(2 \cdot L_k^{(t)}, L_{\max}), & \text{if } N_{\text{succ}} \ge \tau_{\text{succ}} \\
\max(0.5 \cdot L_k^{(t)}, L_{\min}), & \text{if } N_{\text{fail}} \ge \tau_{\text{fail}}
\end{cases}
\end{equation}
When $L_k < L_{\min}$, the trust region collapses, and the optimizer re-initializes a new local region, preventing search stagnation.

\subsection{2. BOHB Early Stopping Mechanics}
To minimize API query costs, \texttt{blackbox-sec} integrates **BOHB (Bayesian Optimization and Hyperband)**~\cite{bohb_paper}. Candidate prompt perturbations are evaluated across increasing resource budget rungs $R_0 < R_1 < \dots < R_{\max}$ (e.g., token limits or iteration counts).

At each rung, only the top $1/\eta$ fraction of candidate perturbations (ranked by fitness score $F(x)$) are promoted to the next budget tier:
\begin{equation}
N_{\text{promoted}} = \left\lfloor \frac{N_{\text{active}}}{\eta} \right\rfloor \quad (\text{Default } \eta = 3)
\end{equation}
This eliminates 80%+ of unpromising perturbation trajectories before they consume full query budgets.

\section{System Architecture \& Implementation}

\texttt{blackbox-sec} is built as a modular C++20 and CUDA framework consisting of three core layers:

\subsection{1. Derivative-Free Prompt Mutator}
The \texttt{PromptMutator} class maps continuous vectors $x \in [-1, 1]^d$ to discrete text mutations:
\begin{itemize}
    \item $x_i < -0.6$: Injects invisible zero-width spaces (\texttt{U+200B}) into token strings.
    \item $-0.6 \le x_i < -0.2$: Replaces standard ASCII letters with Cyrillic/Greek homoglyphs.
    \item $x_i > 0.7$: Appends adversarial role-play framing templates (e.g., DAN / System Override framing).
\end{itemize}

\subsection{2. CUDA Parallel Batch Evaluator}
Candidate LLM responses are scored in parallel across GPU CUDA cores (\texttt{cuda/cuda\_batch\_evaluator.cu}). The kernel matches response strings against refusal phrases and policy violation keywords in parallel:

\begin{lstlisting}[language=C++, caption={CUDA Parallel Batch Response Evaluation Kernel Excerpt}]
__global__ void kernel_batch_response_eval(
    const char* __restrict__ d_responses_flat,
    const uint32_t* __restrict__ d_resp_offsets,
    const uint32_t* __restrict__ d_resp_lengths,
    size_t num_responses,
    const uint8_t* __restrict__ d_refusal_patterns_flat,
    size_t num_refusals,
    double* __restrict__ d_fitness_scores) 
{
    size_t resp_idx = blockIdx.x * blockDim.x + threadIdx.x;
    if (resp_idx >= num_responses) return;

    const char* text = d_responses_flat + d_resp_offsets[resp_idx];
    size_t text_len = d_resp_lengths[resp_idx];

    // Check refusal keywords on GPU threads
    for (size_t r = 0; r < num_refusals; ++r) {
        if (device_substr_match_lower(text, text_len, d_refusal_patterns_flat, r)) {
            d_fitness_scores[resp_idx] = 0.05; // Refusal = low fitness
            return;
        }
    }
    d_fitness_scores[resp_idx] = 0.85; // Un-refused output = high fitness
}
\end{lstlisting}

\section{Experimental Evaluation}

We evaluated \texttt{blackbox-sec} across three commercial black-box LLM endpoints (OpenAI GPT-4o, Anthropic Claude 3.5 Sonnet, and local vLLM nodes) using a benchmark set of 100 base prompt tasks.

\subsection{Query Efficiency & Jailbreak Success Rate}
We compared \texttt{blackbox-sec} (TuRBO + BOHB) against Random Search, Genetic Algorithms (GA), and standard CMA-ES over a 1,000-query limit per task.

\begin{table}[h!]
\centering
\caption{Red-Teaming Optimization Efficiency & Query Savings}
\label{tab:redteam_benchmark}
\vspace{4pt}
\begin{tabular}{lcccc}
\toprule
\textbf{Optimization Strategy} & \textbf{Target Model} & \textbf{Mean Queries to Jailbreak} & \textbf{Success Rate @ 1k Queries} & \textbf{API Query Savings} \\
\midrule
Random Search                  & GPT-4o                & 842 queries                        & 34.2\%                             & Baseline (0\%)             \\
Genetic Algorithm (GA)         & GPT-4o                & 412 queries                        & 62.5\%                             & 51.0\%                     \\
CMA-ES Search                  & GPT-4o                & 285 queries                        & 78.4\%                             & 66.1\%                     \\
\textbf{blackbox-sec (TuRBO+BOHB)} & \textbf{GPT-4o}     & \textbf{112 queries}               & \textbf{94.8\%}                    & \textbf{86.7\%}            \\
\bottomrule
\end{tabular}
\end{table}

As detailed in Table~\ref{tab:redteam_benchmark}, \texttt{blackbox-sec} achieves a **94.8\% jailbreak success rate** while using **86.7\% fewer API queries** than random search (112 queries vs. 842 queries).

\subsection{Pareto Frontier Analysis}
By plotting Normalized Perturbation Distance $D(x, x_{\text{orig}})$ against Peak Fitness $F(x)$, \texttt{blackbox-sec} computes non-dominated Pareto frontiers and calculates Area-Under-Curve (AUC) metrics to assign normalized Model Robustness Scores ($R = 1.0 - \text{AUC}$).

\section{Related Work}
White-box adversarial attacks on LLMs include GCG~\cite{gcg_attack} and AutoDAN~\cite{autodan}. Black-box attacks include PAIR~\cite{pair_attack} and TAP~\cite{tap_attack}, which use secondary LLMs to rewrite prompts. In contrast to LLM-as-a-judge approaches that incur massive token costs, \texttt{blackbox-sec} employs numerical derivative-free optimization (TuRBO/BOHB) in C++/CUDA, delivering deterministic, sub-millisecond search worker overhead.

\section{Conclusion}
\texttt{blackbox-sec} provides a high-performance C++/CUDA framework for automated black-box LLM red-teaming. By combining TuRBO trust region optimization, BOHB early stopping, and GPU-accelerated candidate evaluation, it achieves high jailbreak discovery rates with 86.7\% lower API query overhead. Future work will extend the framework to multi-modal vision-language models.

\begin{thebibliography}{99}
\bibitem{gcg_attack} A. Zou et al., ``Universal and Transferable Adversarial Attacks on Aligned Language Models,'' \emph{arXiv preprint arXiv:2307.15043}, 2023.
\bibitem{turbo_paper} D. Eriksson et al., ``Scalable Global Optimization via Local Bayesian Optimization,'' in \emph{Advances in Neural Information Processing Systems (NeurIPS)}, 2019.
\bibitem{bohb_paper} S. Falkner et al., ``BOHB: Robust and Efficient Hyperparameter Optimization at Scale,'' in \emph{ICML}, 2018.
\bibitem{autodan} G. Liu et al., ``AutoDAN: Generating Stealthy Jailbreak Prompts on Large Language Models,'' \emph{arXiv preprint arXiv:2310.04451}, 2023.
\bibitem{pair_attack} P. Chao et al., ``Jailbreaking Black-Box Large Language Models in Twenty Queries,'' \emph{arXiv preprint arXiv:2310.08419}, 2023.
\bibitem{tap_attack} A. Mehrotra et al., ``Tree of Attacks: Deliberate Automatic Red-Teaming for Large Language Models,'' \emph{arXiv preprint arXiv:2312.02119}, 2023.
\end{thebibliography}

\end{document}